\documentclass[12pt,a4paper]{article}

\usepackage[utf8]{inputenc}
\usepackage[T1]{fontenc}
\usepackage[french]{babel}

\usepackage{styles/couleurs}
\usepackage{styles/mise_en_page}
\usepackage{styles/boites}
\usepackage{styles/maths}
\usepackage{styles/figures}

\usepackage{amsmath}
\usepackage{amssymb}
\usepackage{enumitem}
\usepackage{array}
\usepackage{tabularx}
\usepackage{graphicx}
\usepackage{tikz}
\usepackage{xparse}
\usetikzlibrary{calc,angles,quotes}

\setlist[itemize]{leftmargin=1.4em,itemsep=0.15em,topsep=0.2em}
\setlist[enumerate]{leftmargin=*,itemsep=0.22em,topsep=0.22em}

\newcommand{\titresequence}[1]{%
\begin{center}
  {\sffamily\LARGE\bfseries\color{TitreC}Séquence 2 -- Théorème de Pythagore}\\[0.2em]
  {\sffamily\large\bfseries\color{SavoirC}#1}
\end{center}}

\newcommand{\lignepointillee}{\noindent\dotfill\par\vspace{0.45em}}
\newcommand{\lignesreponse}[1]{%
\begin{tcolorbox}[
  enhanced,colback=white,colframe=NeutreC!40,boxrule=0.5pt,arc=1.5mm,
  left=3mm,right=3mm,top=2mm,bottom=1mm,before skip=2pt,after skip=2pt]
  \foreach \i in {1,...,#1}{\lignepointillee}
\end{tcolorbox}}

\begin{document}
\titresequence{Automatismes -- Séries courtes}

\remarque{Chaque série est prévue pour environ cinq minutes.}

\section*{Série 1 -- Hypoténuse}
Dans chaque cas, nomme l'hypoténuse.
\begin{enumerate}[label=\alph*)]
  \item $ABC$ rectangle en $B$ ;
  \item $DEF$ rectangle en $F$ ;
  \item $MNP$ rectangle en $M$ ;
  \item $RST$ rectangle en $S$.
\end{enumerate}

\section*{Série 2 -- Égalité de Pythagore}
Écris l'égalité adaptée à chaque triangle précédent.

\section*{Série 3 -- Carrés}
\[
3^2=\dotfill\quad 4^2=\dotfill\quad 5^2=\dotfill\quad
8^2=\dotfill\quad 10^2=\dotfill\quad 12^2=\dotfill
\]

\section*{Série 4 -- Racines carrées simples}
\[
\sqrt{9}=\dotfill\quad \sqrt{16}=\dotfill\quad \sqrt{25}=\dotfill\quad
\sqrt{49}=\dotfill\quad \sqrt{64}=\dotfill\quad \sqrt{100}=\dotfill
\]

\section*{Série 5 -- Addition ou soustraction ?}
Indique l'opération finale utile, sans faire le calcul.
\begin{enumerate}[label=\alph*)]
  \item On cherche l'hypoténuse et on connaît les deux autres côtés.
  \item On cherche un côté de l'angle droit et on connaît l'hypoténuse.
\end{enumerate}

\section*{Série 6 -- Triplets}
Complète.
\[
3-4-\dotfill
\qquad
5-12-\dotfill
\qquad
6-8-\dotfill
\qquad
\dotfill-15-17
\]

\section*{Série 7 -- Vrai ou faux ?}
\begin{enumerate}
  \item L'hypoténuse est toujours le plus grand côté.
  \item Le théorème de Pythagore s'utilise dans tous les triangles.
  \item Si $AB^2=25$, alors $AB=5$.
  \item $\sqrt{36}=18$.
\end{enumerate}
\end{document}
